\documentclass[../main.tex]{subfiles}
\begin{document}
%%%%%%%%%%%%%%%%%%%%%%%%%%%%%%%%%%%%%%%%%%%%%%%%%%%%%%%%%%%%%%%%%%%%%%%%%%%%%%%%%%%%%%%%%%%%%%%%%%%
\chapter{Getting Synced with OmniPresence and Omni Sync Server}

Available on the Mac and iOS, OmniGraffle is an essential tool for all your design and
diagramming needs. And when it comes to protecting those documents and sharing them between
your Mac and iOS devices, we have something a little extra-special for you:
\textbf{OmniPresence} and \textbf{Omni Sync Server}.

\begin{itemize}
  \item \textbf{OmniPresence} is a free utility that you install on your Mac for synchronizing
    files. OmniPresence runs in the background, watching a ``connected folder.'' It makes sure
    the files inside are in sync with files in the cloud (a WebDAV server of your choice), and
    with your OmniPresence-enabled iPad apps.

    On iOS, OmniPresence is built into all of our current document-based apps, including
    OmniGraffle 2, OmniOutliner 2, and OmniPlan 3. You'll access OmniPresence folders and files
    from within the app's respective Resource Browsers. The result is that all of your
    OmniPresence-enabled devices detect changes wherever you make them, and constantly sync up
    to keep your files up to date.

    To get started, you'll need to
    \href{http://www.omnigroup.com/download/latest/omnipresence}{download OmniPresence for Mac}
    from our website.
\end{itemize}

\begin{ognote}
  Unlike other Omni apps, OmniPresence isn't a replacement for your current sync solution for
  OmniFocus. OmniFocus uses its own built-in sync code (rather than OmniPresence) to push
  database changes to a WebDAV server and retrieve them from it. OmniFocus for Mac doesn't
  require OmniPresence to be installed for its sync to function to work.
\end{ognote}

\begin{itemize}
  \item \textbf{Omni Sync Server} resides in a well-air conditioned secure facility. Omni Sync
    Server is our own custom WebDAV server and it functions as a central hub for all your
    syncing needs; it's what OmniPresence connects to when shuttling files between your Mac and
    iOS devices.

    With Omni Sync Server, we never ever look at your data without your permission. For example,
    when our Support Humans are working with you to restore a corrupt file or isolate a
    particular issue you're having with our server. Beyond that, you can rest assured that your
    data is safe and secure on Omni Sync Server.
\end{itemize}

\begin{ogsidebar}
  Keeping Your Data Private with OmniPresence

  We understand that sometimes you want to keep your data on your own server. This could be due
  to company policy, for legal reasons, or just your desire to be in charge of backing up
  everything that's important to you. We get it!

  If you don't want to use Omni Sync Server, OmniPresence is capable of connecting to any WebDAV
  server you choose. This could be one that you run on your own server, or a third-party
  service that you rely on.

  If you choose to run your own WebDAV server, step-by-step instructions for setting up WebDAV
  on macOS Server are available
  \href{http://support.omnigroup.com/omnipresence-os-x-server}{in this handy support article
  on our website}.
\end{ogsidebar}

\section{Creating an Omni Sync Server Account}

Your first order of business---if you are planning to use OmniPresence with Omni Sync
Server---is to create an Omni Sync Server account. In doing so, you will create a Account Name
and password, which you will use when connecting OmniPresence on your Mac.

Follow along\ldots{}

\begin{enumerate}
  \item Using your favorite web browser, and go to the
    \href{https://manage.sync.omnigroup.com}{Omni Sync Server webpage}:

    \ogfigure{0.800}{res/op14mac_osscreate01.png}
  \item At the top of the page, click \textbf{Sign Up}.
  \item In the fields provided, enter the following information:

    \ogfigure{0.800}{res/op14mac_osscreate02.png}

    \begin{itemize}
      \item \textbf{Account Name} --- You will use your Account Name when connecting to Omni Sync
        Server. Account Names can include lowercase letters, numbers, or symbols---sorry,
        uppercase letters, spaces, and the following symbols are not permitted:
        \textbf{\textasciitilde{} ` : /} and \textbf{\textbackslash\{\}}.
      \item \textbf{Email Address} --- Enter your email address.
      \item \textbf{Password} --- Enter a unique, and secure, password. After entering your
        password, you'll need to enter it again in the next field, just to verify that you have
        entered the password correctly.
    \end{itemize}
  \item \textbf{Terms of Use} --- Please take some time to read through the Omni Sync Server
    Terms of Use. If you agree with the terms and conditions, click the checkbox.
  \item Finally, click \textbf{Create Account}.
\end{enumerate}

After clicking the \textbf{Create Account} button, we will send a confirmation message to the
email address you specified. Keep an eye out for that email; when you receive it, be sure to
click the \textbf{Finish Registration} button in that email to activate your Omni Sync Server
account.

\begin{ognote}
  Keep the Account Name and Password used to create your Omni Sync Server account handy, as you
  will need them in the next chapter when you install and configure OmniPresence on your Mac.
\end{ognote}

\section{Installing and Configuring OmniPresence}

OmniPresence starts out its life on your Mac just like any other Mac app you install. You
download a \textbf{disk image}, and drag the app icon to the Applications folder. However,
unlike other Mac apps that show up in the Dock when you launch them, OmniPresence never shows
up in your Dock. Instead, OmniPresence spends its life as a \textbf{menu bar extra}---one of
those black icons---at the top of the screen.

With the Account Name and Password for your Omni Sync Server account in-hand, this chapter takes
you through the process of downloading, installing, and configuring OmniPresence on your Mac.

\subsection{Downloading OmniPresence from our Server}

OmniPresence is only available for download from our website; it isn't available on the Mac App
Store.

Click \href{http://www.omnigroup.com/download/latest/omnipresence}{this link to download
OmniPresence}. OmniPresence is packaged within a disk image (a file ending with a
\textit{.dmg} file extension), which is saved to the \textbf{Downloads} folder on your Mac.

\ogfigure{0.800}{res/op14mac_sync01.png}

\subsection{Installing OmniPresence}

To install OmniPresence, follow these steps:

\begin{enumerate}
  \item Double-click \textbf{OmniPresence-1.4.1.dmg} to open the disk image. This opens the disk
    image, mounting it under the \textbf{Devices} section of the Finder's sidebar.

    \begin{ognote}
      The filename you see might be slightly different. As of this writing, OmniPresence for Mac is
      at version 1.4.1 (thus the filename, \textit{OmniPresence-1.4.1.dmg}). If you see a
      different number in the filename, don't be worried; just go ahead and double-click that one.
      It just means that our engineering team have been hard at work to improve OmniPresence for
      Mac.
    \end{ognote}
  \item After agreeing to our license, you're presented with the following window:

    \ogfigure{0.800}{res/op14mac_sync02.png}
  \item Drag and drop the \textbf{OmniPresence.app} icon onto the Applications folder to install
    OmniPresence on your Mac:

    \ogfigure{0.800}{res/op14mac_sync03.png}
  \item Close that window by pressing \textbf{Command-W}.
  \item In the Finder's sidebar, click the eject button next to OmniPresence's disk image; this
    unmounts the disk image so you can delete the disk image file from your system.

    \ogfigure{0.800}{res/op14mac_ejectfmfinder.png}
  \item In the Finder window, return to the \textbf{Downloads} folder. Select the file
    \textit{OmniPresence-1.4.1.dmg} and press \textbf{Command-Delete} to move the file to the
    Trash. The next time you empty the Trash, that disk image is purged from your system, saving
    you precious disk space.
\end{enumerate}

With OmniPresence installed, now it's time to set up a synced folder for OmniPresence to sync
to, and then connect to Omni Sync Server.

\subsection{Starting OmniPresence and Connecting to Omni Sync Server}

With OmniPresence installed, it's time to get it started and connect to your Omni Sync Server
account.

Follow along\ldots{}

\begin{enumerate}
  \item In the Finder window, choose \textbf{Applications} in the Finder's sidebar. (If you
      don't see it there, choose \textbf{Go} $\blacktriangleright$ \textbf{Applications}, or use
    \textbf{Shift-Command-A}.)

    \ogfigure{0.800}{res/op14mac_sync04.png}
  \item Scroll down in the list of applications installed on your Mac and locate
    \textbf{OmniPresence.app}.
  \item Double-click the OmniPresence.app icon to open the app. You'll notice that, unlike other
    Mac apps, OmniPresence's icon doesn't appear in the Dock. Instead, OmniPresence installs a
    \textit{menu bar extra} --- \ogicon{11.1pt}{res/op14mac_icon_syncready.png} --- in the menu
    bar at the top of the screen.

    \ogfigure{0.800}{res/op14mac_sync04a.png}

    For more details about using OmniPresence's menu bar extra---and the icon's various
    states---see Using OmniPresence on Your Mac.
  \item Click \textbf{Set Up Syncing}:

    \ogfigure{0.800}{res/op14mac_sync05.png}
  \item Click \textbf{Create a Synced Folder}:

    \ogfigure{0.800}{res/op14mac_sync06.png}
  \item In the sheet that appears, you have the option to name the OmniPresence folder that
    lives on your Mac. We've pre-filled the \textbf{New Folder} field with
    \textbf{OmniPresence}, but you can change this to whatever you'd like.

    You have total control over where the sync folder is stored on your Mac. For example, you
    could create it within the \textbf{Documents} folder, or you could opt to have your
    OmniPresence folder reside at the bottom level of your Home folder.

    \ogfigure{0.800}{res/op14mac_sync07.png}

    \ogfigure{0.800}{res/op14mac_sync07a.png}

    When you've settled on a name and location for your OmniPresence folder---or if you've decided
    to go with the default options---click \textbf{Save}. This creates the folder OmniPresence
    uses to synchronize your files.
  \item After creating the OmniPresence folder, your next task is to choose a server in the
    cloud where OmniPresence stores your synced files.

    You have two options to choose from:

    \ogfigure{0.800}{res/op14mac_sync08.png}

    \begin{itemize}
      \item \textbf{Omni Sync Server}: This is Omni's free file server, which has specifically
        been designed to work in unison with OmniPresence.
      \item \textbf{Web Server}: This could be your own WebDAV server, or a WebDAV server hosted
        by a third-party provider.
    \end{itemize}
  \item If you have an Omni Sync Server account, enter your \textbf{Account Name} and
    \textbf{Password} in the applicable fields and click \textbf{Connect}.

    \begin{ognote}
      If you haven't already set up an Omni Sync Server account, click \textbf{Sign Up for an Omni
      Sync Server Account}. This opens your default web browser and transports you to the Omni
      Sync Server site where you can sign up for an account and set your password.
    \end{ognote}

    After clicking \textbf{Connect}, OmniPresence connects to the server you've chosen and tests
    for server compatibility. This is more of an issue with third-party WebDAV servers than it
    is for connecting to Omni Sync Server, which has been designed to work with OmniPresence.

    If everything is okay, OmniPresence makes the connection and lets you know if it is able to
    connect to Omni Sync Server (or the WebDAV server you've specified).

    \ogfigure{0.800}{res/op14mac_sync09.png}
  \item If you would like OmniPresence to start up whenever you log into your Mac, click the
    checkbox for \textbf{Open OmniPresence at Login}.

    \ogfigure{0.800}{res/op14mac_sync10.png}
  \item At the bottom of that window, click \textbf{Done} if you're finished, or click
    \textbf{Open Synced Folder} to open a Finder window displaying your new sync folder.

    If you've already synced with OmniPresence on another device, OmniPresence automatically
    starts populating this folder with your existing synced files. Otherwise, add files to this
    folder to start syncing them with the server.

    \ogfigure{0.800}{res/op14mac_sync11.png}
\end{enumerate}

Once OmniPresence connects with Omni Sync Server, you're ready to go. Whenever you create files
with OmniGraffle, OmniOutliner, or OmniPlan, be sure to save them to your OmniPresence folder.
And, if you have OmniGraffle, OmniOutliner, or OmniPlan for iOS, don't forget to connect to
Omni Sync Server on your device so your files will be with you wherever you go.

\section{Using OmniPresence on Your Mac}

OmniPresence opens up a world of possibilities for syncing---and sharing---files between your
Mac and iOS devices. On iOS, you access your synced folders from the Document Browser in
OmniGraffle, OmniOutliner, or OmniPlan. On the Mac, however, you use the Finder for file
management.

This final section discusses ways to \textit{use} OmniPresence on your Mac, and from within
OmniGraffle, OmniOutliner, and OmniPlan to help streamline your workflow. You'll also learn
some valuable tips about how to deal with conflict files that might occur when the same synced
file is accessed on two different Macs or devices.

\subsection{How Much Data Can I Store on Omni Sync Server?}

While we don't have a hard limit, we ask that you limit your data usage to 1 GB or less. This
ensures that all our customers are able to get the best performance from their free Omni Sync
Server account.

To determine how much data is stored in your OmniPresence folder, select the folder in the
Finder and press \textbf{Command-I}. This opens the \textbf{Get Info} window, which not only
displays how much data is stored in the folder, but also the number of files within, too.

\ogfigure{0.800}{res/op14mac_getinfowindow.png}

If you are using more than 1 GB, we may ask you to reduce the amount of data you are storing. It
is possible that we will implement a limit in the future, so it's a good idea to keep your
usage of the Omni Sync Server to a reasonable size.

If your workflow requires more storage than 1 GB, you may want to either host your own server,
or use a third-party paid WebDAV host. In that case, the only limits are the ones you set for
yourself.

\begin{ogtip}
  If you are interesting in hosting your own server, instructions can be found on our support
  site, in \href{https://support.omnigroup.com/omnipresence-os-x-server}{Configure WebDAV for
  OmniPresence}.
\end{ogtip}

\section{Getting to Know the Menu Bar Status Icons}

As noted earlier, the way you interact with OmniPresence as an app is not by clicking an app
icon in the Dock, but by clicking its icon in the menu bar.

\begin{description}[style=nextline]
  \item[{\ogicon{11.1pt}{res/op14mac_icon_syncready.png}}] In its most basic state, OmniPresence
    is online and ready to sync files.
  \item[{\ogicon{11.1pt}{res/op14mac_icon_syncinprogress.png}}] Whenever OmniPresence notices a
    change to the files in the sync folder, the icon appears to be lifting a barbell.
  \item[{\ogicon{11.1pt}{res/op14mac_icon_syncoffline.png}}] If you lose your Internet
    connection (which is required for syncing), the icon turns gray to denote that you are
    offline.
  \item[{\ogicon{11.1pt}{res/op14mac_icon_syncerror.png}}] If you see a gray sync icon with an
    exclamation mark, it means that OmniPresence failed to connect to the sync server. This
    could be the result of an invalid Account Name or Password, or something as simple as an
    intermittent Internet connection that drops while OmniPresence is trying to connect to the
    sync server.

    If the icon remains in this state, here are some things you could try:

    \begin{itemize}
      \item Double-check your \textbf{Account Name} and \textbf{Password}.
      \item Check your network connection in \textbf{System Preferences} $\blacktriangleright$
        \textbf{Network}.
      \item Try quitting and restarting OmniPresence.
    \end{itemize}

    If that doesn't help, contact our Support Humans for assistance.
  \item[{\ogicon{11.1pt}{res/op14mac_icon_syncpaused.png}}] If you pause syncing by clicking the
    \textbf{Pause} button in the popover menu, the icon turns gray and gains a pause symbol. To
    resume syncing, click the icon to open the popover menu, and then click
    \ogicon{10.4pt}{res/op14mac_icon_syncresume.png}. If you are connected to the Internet,
    OmniPresence springs back to life and syncs with the server.
\end{description}

\section{Creating Subfolders in Your Sync Folder}

After creating your OmniPresence sync folder, there's nothing that's stopping you from just
dumping all of your OmniGraffle, OmniOutliner, or OmniPlan files in there and calling it a
day. But why be so cavalier when you can create subfolders for organizing your files?

Follow along\ldots{}

\begin{enumerate}
  \item Open a Finder window and select your OmniPresence folder:

    \ogfigure{0.800}{res/op14mac_subfolder01.png}
  \item Choose \textbf{File} $\blacktriangleright$ \textbf{New Folder}, or use the keyboard
    shortcut, \textbf{Shift-Command-N} to add an \textbf{untitled folder} inside the
    OmniPresence folder:

    \ogfigure{0.800}{res/op14mac_subfolder02.png}
  \item Enter a new name for the folder, such as \textbf{Graffle Files}, or \textbf{Outlines},
    or \textbf{Science Project}---whatever the case may be:

    \ogfigure{0.800}{res/op14mac_subfolder03.png}
  \item After renaming the folder, press \textbf{Return} to accept the name:

    \ogfigure{0.800}{res/op14mac_subfolder04.png}
\end{enumerate}

\begin{ognote}
  OmniPresence \textit{does not} sync empty folders. Instead, OmniPresence waits until something
  is stored in the folder before the file \textit{and the folder} are synced to the server.
\end{ognote}

There is no limit (except for the 1 GB storage limit) to the number of subfolders and
sub-subfolders your OmniPresence folder can contain. If you delete a folder, OmniPresence
recognizes that, too, and will remove that folder (and any files it contains) from Omni Sync
Server.

\section{OmniPresence Tips}

Here are some helpful tips for enhancing your OmniPresence experience:

\begin{itemize}
  \item Click the Action menu (the gear next to the Pause/Resume button) for a host of
    OmniPresence options, including:

    \ogfigure{0.800}{res/op_sync12.png}

    \begin{itemize}
      \item About OmniPresence
      \item Release Notes
      \item Open at Login
      \item Check for Updates
      \item Preferences and Privacy
      \item OmniPresence Help
      \item Contact Omni
      \item Quit OmniPresence
    \end{itemize}
  \item To stop syncing a folder you've connected, click its \textbf{Info} button and choose
    \textbf{Disconnect}. The folder on your Mac is moved to the Trash and the files inside will
    no longer be synced. Any files left in the folder before disconnecting will stay safe and
    sound at the cloud server location you've specified. To reconnect and start syncing with
    those files again, just create a new folder connected to that cloud location.

    \ogfigure{0.800}{res/op_sync14.png}
  \item Unsure of when OmniPresence last synced? Click the OmniPresence icon in the menubar and
    then look beneath each folder name to see the date and time that folder last synced.

    \ogfigure{0.800}{res/op_sync13.png}
\end{itemize}

\begin{ognote}
  If you don't set up sync now, but later decide that you want to, you can always install
  OmniPresence and set up an Omni Sync Server account when it's more convenient for you.
\end{ognote}

\section{Linking Template and Stencil Folders on your Mac to OmniGraffle}

To make it easier for you to quickly find your OmniGraffle stencils and template files, you can
create a link to those folders in the Resource Browser. You can link to folders on your Mac's
hard drive, as well as to folders you are syncing, such as those with OmniPresence or on
iCloud Drive.

Using the Resource Browser (\textbf{File} $\blacktriangleright$ \textbf{Resource Browser}), you
can add a \textbf{Linked Folder} to the \textbf{Templates} and \textbf{Stencils} sections of
the sidebar.

\begin{enumerate}
  \item In the Resource Browser, click \ogicon{12.5pt}{res/og7_button_linkedfolder.png}:

    \ogfigure{0.800}{res/og7_sync_linkedfolder01.png}
  \item Choose a folder in the sheet that appears and click \textbf{Open} to add the folder to
    the Resource Browser's sidebar:

    \ogfigure{0.800}{res/og7_sync_linkedfolder02.png}
\end{enumerate}

After you click \textbf{Open}, OmniGraffle examines the contents of the linked folder, and
places a link to that folder to both the Templates and Stencils sections of the Resource
Browser's sidebar.

\ogfigure{0.800}{res/og7_sync_linkedfolder03.png}

To remove a Linked Folder from the Resource Browser's sidebar, first select the folder in
question, click \ogicon{12.5pt}{res/og7_button_rb_actionmenu.png}, and then choose
\textbf{Remove Linked Folder}.
%%%%%%%%%%%%%%%%%%%%%%%%%%%%%%%%%%%%%%%%%%%%%%%%%%%%%%%%%%%%%%%%%%%%%%%%%%%%%%%%%%%%%%%%%%%%%%%%%%%
\end{document}
